\documentclass[12pt]{article}
\usepackage[utf8]{inputenc}
\usepackage[T1]{fontenc}
\usepackage{hyperref}
\title{Tai Pao Language Resources for K-12 Curriculum Development}
\author{Compiled by Claude Sonnet 5 (claude-sonnet-5) for Andr\'as Kornai}
\date{}
\begin{document}
\maketitle
\section{Introduction}
Tai Pao (ISO 639-3 tpo, Glottocode taip1250) is a Tai language spoken in Bolikhamxai Province, Laos (Khamkeut, Pakkading and Viengthong districts), and in Vietnam, where it is called Tai Hang Tong. It is closely related to Tai Yo. Speaker numbers vary between sources (roughly 13,000 in older estimates). This report is based on a single round of web searches. No educational or linguistic resources specific to Tai Pao were found.
\section{Primary Educational Resources}
None found.
\section{Linguistic and Reference Materials}
General reference pages only:

\url{https://en.wikipedia.org/wiki/Tai_Pao_language}

\url{https://glottolog.org/resource/languoid/iso/tpo}
\section{Religious and Community Resources}
The Joshua Project language page notes that Bible translation has started, but it lists no available scripture:

\url{https://joshuaproject.net/index.php/languages/tpo}

Asia Harvest people-group profile for Laos: \url{https://people-groups.asiaharvest.org/Laos/TaiPao-PBW.pdf}
\section{Digital and Audio Resources}
None found.
\section{Curriculum Development Context}
The language is not known to be taught in schools, and no materials were found. The search results mention a distinct local script in Nghe An, Vietnam (Lai Pao), but no resources for it were located. Any curriculum would have to begin with new documentation.
\end{document}